\scanpage{77}
\mspage{76}

alors le $1^{\mathrm{er}}$ \unsure{mb|mbr.} de (*) \deleted{\unsure{s'écrit}} équivaut à
\[
A\subset B
\]
qui implique bien $|A|\subset|B|$ tautologiquement. Ainsi

\uline{Lemme \unsure{!|1}} Sous la \inserted{seule} condition $(\mathrm{iii}_M)$, si $A$ et $B$ sont des parties de $M$, on a
\diagram{
\begin{tikzcd}[ampersand replacement=\&,column sep=large,row sep=small]
\operatorname{supp}A\subset\operatorname{supp}B
\arrow[d,phantom,"{\unsure{\Vert|\Updownarrow}}" description]
\& |A|\subset|B|\arrow[l,Rightarrow]\\
(A\subset\operatorname{supp}B)\&
\end{tikzcd}
}
et l'implication inverse est vraie si $B$ est un support (i.e.\ $B=\operatorname{supp}B$).

\uline{Corollaire} L'application
\[
\Sigma_M\longrightarrow\unsure{\mathcal{P}|\mathscr{P}}(P)\qquad A\longmapsto|A|
\]
\begin{tabular}{|p{0.88\linewidth}}
est injective, et la relation d'inclusion dans $\Sigma_M$ est induite par la relation d'inclusion dans $\mathcal{P}(P)$
\end{tabular}

Qu'en est-il de l'intersection $\bigwedge_i$ \inserted{$=\bigcap_i$} et du \deleted{la ``réunion''} sup $\bigvee_i$ des supports?

On a
\[
\begin{aligned}
\left|\bigwedge_i A_i\right|&\subset\bigcap_i|A_i|&\text{tautologique}\\
\left|\bigvee_i A_i\right|&\supset\bigcup_i|A_i|&\text{id}
\end{aligned}
\]
mais je doute qu'on ait tjrs égalité. D'autre part, je dis que
\[
|\unsure{c_M|C_M}(A)|\subset\unsure{R|P}\setminus|A|
\]
